\chapter{Presupuesto}

El presupuesto de la obra se elabora a partir del inventario de 116 puntos de red definido en el capítulo de Planos, considerando los materiales pasivos, los equipos activos, la mano de obra y la certificación final. Todos los montos son \textbf{referenciales} y están expresados en soles (S/) con IGV incluido; corresponden a precios promedio del mercado nacional de telecomunicaciones y pueden variar según el proveedor y la fecha de compra.

\section{Presupuesto General de la Obra}
El presupuesto general agrupa la inversión en tres grandes rubros: los materiales y equipos del sistema de cableado, la mano de obra de instalación y la certificación de los enlaces. El resumen se muestra en el Cuadro~\ref{tab:presupuesto_general}.

\begin{table}[H]
\centering
\caption{Presupuesto general de la obra.}
\label{tab:presupuesto_general}
\begin{tabular}{clr}
\hline
\textbf{Ítem} & \textbf{Descripción del Rubro} & \textbf{Monto (S/)} \\
\hline
1 & Materiales y equipos (cableado, gabinete, activos) & 34\,094.00 \\
2 & Mano de obra de instalación (estimada en 30\,\%) & 10\,228.00 \\
3 & Certificación de los 116 enlaces (global) & 2\,900.00 \\
\hline
\multicolumn{2}{r}{\textbf{TOTAL ESTIMADO DE LA OBRA}} & \textbf{47\,222.00} \\
\hline
\end{tabular}
\end{table}

\noindent El costo aproximado por punto de red instalado y certificado es de S/ 47\,222.00 / 116 $\approx$ \textbf{S/ 407.00 por nodo}, valor coherente con proyectos de cableado estructurado de mediana escala en el ámbito educativo.

\section{Presupuesto Detallado por Partidas}
A continuación se desglosa el presupuesto en partidas individuales, separando los componentes pasivos (cables, conectores, canalizaciones) de los equipos activos y de respaldo. Esta clasificación facilita la elaboración del cronograma de compras.

\begin{table}[H]
\centering
\caption{Presupuesto detallado por partidas.}
\label{tab:presupuesto_detallado}
\footnotesize
\begin{tabular}{clccrr}
\hline
\textbf{N.\textsuperscript{o}} & \textbf{Descripción de la Partida} & \textbf{Und.} & \textbf{Cant.} & \textbf{P. Unit. (S/)} & \textbf{Parcial (S/)} \\
\hline
\multicolumn{6}{l}{\textit{A. Componentes pasivos (cableado y conectividad)}} \\
%
1  & Cable U/UTP Cat 6A, caja x 305\,m            & caja & 15  & 950.00  & 14\,250.00 \\
2  & Patch Panel 24 puertos Cat 6A               & und  & 5   & 200.00  & 1\,000.00 \\
3  & Jack Keystone RJ-45 Cat 6A                  & und  & 125 & 18.00   & 2\,250.00 \\
4  & Faceplate (placa de pared)                  & und  & 107 & 8.00    & 856.00 \\
5  & Caja de paso/registro (APs)                 & und  & 9   & 12.00   & 108.00 \\
6  & Patch cord Cat 6A 1\,m (parcheo)            & und  & 116 & 15.00   & 1\,740.00 \\
7  & Organizador horizontal de cable 1U          & und  & 6   & 60.00   & 360.00 \\
8  & Canaleta PVC 40$\times$25\,mm (tramo 2\,m)  & und  & 220 & 14.00   & 3\,080.00 \\
9  & Tubería conduit y accesorios                & glb  & 1   & 1\,000.00 & 1\,000.00 \\
%
\multicolumn{6}{l}{\textit{B. Equipos activos y de respaldo}} \\
%
10 & Gabinete de piso 42U con llave              & und  & 1   & 1\,800.00 & 1\,800.00 \\
11 & Switch 48 puertos Gigabit PoE+              & und  & 3   & 1\,900.00 & 5\,700.00 \\
12 & UPS Online 1500\,VA / 1350\,W               & und  & 1   & 1\,500.00 & 1\,500.00 \\
%
\multicolumn{6}{l}{\textit{C. Puesta a tierra y rotulado}} \\
%
13 & Barra de tierra (TGB) y cable verde         & glb  & 1   & 250.00  & 250.00 \\
14 & Etiquetas y material de rotulado            & glb  & 1   & 200.00  & 200.00 \\
\hline
\multicolumn{5}{r}{\textbf{SUBTOTAL MATERIALES Y EQUIPOS}} & \textbf{34\,094.00} \\
\hline
\end{tabular}
\end{table}

\section{Metrados de Materiales}
El metrado consiste en calcular las cantidades físicas de cada material a partir del diseño. El insumo más importante es el cable de cobre, cuyo metrado se obtiene multiplicando el número de enlaces por la longitud promedio de cada uno y añadiendo un margen de desperdicio.

\begin{table}[H]
\centering
\caption{Metrado del cable de cobre horizontal.}
\label{tab:metrado_cable}
\begin{tabular}{lr}
\hline
\textbf{Parámetro de Cálculo} & \textbf{Valor} \\
\hline
Número total de enlaces (nodos)                 & 116 enlaces \\
Longitud promedio estimada por enlace           & 35\,m \\
Longitud total de cable (116 $\times$ 35\,m)    & 4\,060\,m \\
Margen de desperdicio y reservas (10\,\%)       & 406\,m \\
\hline
\textbf{Total de cable requerido}               & \textbf{4\,466\,m} \\
Presentación comercial (caja x 305\,m)          & 305\,m/caja \\
\textbf{Cajas a adquirir (4\,466 / 305 $\approx$ 14.6)} & \textbf{15 cajas} \\
\hline
\end{tabular}
\end{table}

\noindent El resto de materiales se metra directamente a partir del inventario de puntos: 125 jacks (116 instalados más repuestos), 107 faceplates para los puntos de datos fijos, 9 cajas de paso para los Access Points y 5 patch panels de 24 puertos para alojar los 116 nodos con holgura. La longitud de canaleta se estima en función del recorrido de los muros y de los tramos verticales hacia el Rack G09.

\section{Análisis de Precios Unitarios (APU)}
El Análisis de Precios Unitarios descompone el costo de instalar \textbf{un} punto de red, sumando el material directo, la mano de obra y una pequeña proporción de herramientas. Sirve para sustentar el precio de la partida principal del proyecto. El Cuadro~\ref{tab:apu} muestra el APU del punto de datos típico.

\begin{table}[H]
\centering
\caption{Análisis de Precios Unitarios para un punto de red instalado.}
\label{tab:apu}
\begin{tabular}{lccr}
\hline
\textbf{Recurso} & \textbf{Und.} & \textbf{Cant.} & \textbf{Parcial (S/)} \\
\hline
\multicolumn{4}{l}{\textit{Materiales}} \\
Cable Cat 6A (35\,m a S/ 3.19/m)        & m   & 35   & 111.65 \\
Jack Keystone Cat 6A                    & und & 1    & 18.00 \\
Faceplate                              & und & 0.7  & 5.60 \\
Canaleta y accesorios (prorrateo)       & glb & 1    & 25.00 \\
%
\multicolumn{4}{l}{\textit{Mano de obra}} \\
Técnico instalador (tendido + ponchado) & h   & 1.5  & 30.00 \\
%
\multicolumn{4}{l}{\textit{Herramientas}} \\
Desgaste de herramienta (3\,\% M.O.)    & glb & 1    & 0.90 \\
\hline
\multicolumn{3}{r}{\textbf{COSTO POR PUNTO DE RED}} & \textbf{191.15} \\
\hline
\end{tabular}
\end{table}

\noindent Este valor unitario es referencial y no incluye los equipos activos centralizados (gabinete, switches y UPS), que se presupuestan de forma global por ser únicos para todo el edificio.

\section{Fórmulas Polinómicas}
En obras formales, la fórmula polinómica es una expresión matemática que permite reajustar el presupuesto cuando los precios de los insumos varían en el tiempo (por inflación o cambios de mercado). Agrupa los costos en componentes y asigna a cada uno un coeficiente de incidencia que refleja su peso dentro del total.

Para el presente proyecto, de carácter académico, la estructura de costos se puede representar de manera simplificada con tres componentes principales:

\[
K = a\,\frac{C_c}{C_{c0}} + b\,\frac{E_a}{E_{a0}} + c\,\frac{M_o}{M_{o0}}
\]

\noindent donde:
\begin{itemize}[label=$\bullet$,leftmargin=1.4em]
    \item $K$ es el coeficiente de reajuste del presupuesto.
    \item $C_c$ representa el costo del cable y materiales de cobre (incidencia $a \approx 0.45$).
    \item $E_a$ representa el costo de los equipos activos y el gabinete (incidencia $b \approx 0.35$).
    \item $M_o$ representa el costo de la mano de obra (incidencia $c \approx 0.20$).
    \item El subíndice $0$ indica el precio en la fecha base del presupuesto.
\end{itemize}

\noindent La suma de los coeficientes cumple $a + b + c = 1$. En un proyecto universitario basta con \textbf{presentar la fórmula y justificar las incidencias}; no es necesario aplicar índices oficiales de reajuste.
